\DocumentMetadata{
  pdfversion=2.0,
  pdfstandard=ua-2,
  lang=en-US,
  tagging=on,
  tagging-setup={math/setup=mathml-SE,tabsorder=structure}
}
\documentclass[11pt,letterpaper]{article}
\usepackage[margin=0.68in,includefoot]{geometry}
\usepackage{newcomputermodern}
\usepackage{amsmath,amscd,mathtools}
\providecommand{\square}{\mdlgwhtsquare}
\usepackage[hidelinks]{hyperref}
\hypersetup{
  pdftitle={Algebra PhD Exam, September 2009},
  pdfauthor={University of Florida Department of Mathematics},
  pdfsubject={Graduate mathematics examination},
  pdfdisplaydoctitle=true,
  bookmarksopen=true
}
\setcounter{secnumdepth}{0}
\setlength{\parindent}{0pt}
\setlength{\parskip}{0.28em}
\setlength{\emergencystretch}{3em}
\setlength{\leftmargini}{2.15em}
\setlength{\leftmarginii}{2.65em}
\setlength{\leftmarginiii}{3em}
\renewcommand{\labelenumi}{\textbf{\theenumi.}}
\pagestyle{empty}
\raggedbottom
\allowdisplaybreaks
\newenvironment{examproblems}[1]{%
  \begin{enumerate}%
  \setcounter{enumi}{#1}%
  \setlength{\topsep}{0.35em}%
  \setlength{\partopsep}{0pt}%
  \setlength{\itemsep}{0.48em}%
  \setlength{\parsep}{0.18em}%
}{\end{enumerate}}
\newenvironment{examsubparts}{%
  \begin{enumerate}%
  \setlength{\topsep}{0.18em}%
  \setlength{\partopsep}{0pt}%
  \setlength{\itemsep}{0.16em}%
  \setlength{\parsep}{0.1em}%
}{\end{enumerate}}
\newenvironment{examinstructions}{%
  \par\begingroup\itshape
}{%
  \par\endgroup\vspace{0.18em}
}
\makeatletter
\renewcommand\section{\@startsection{section}{1}{0pt}%
  {-1.25ex plus -.25ex minus -.1ex}%
  {0.55ex plus .1ex}%
  {\normalfont\large\bfseries}}
\renewcommand{\@maketitle}{%
  \newpage\null\vskip -1.2em%
  {\footnotesize\bfseries
    \href{https://www.ufl.edu/}{University of Florida}, \href{https://math.ufl.edu/}{Department of Mathematics}\par}%
  \vskip 0.18em%
  \tagstructbegin{tag=Title}%
    {\LARGE\bfseries\href{https://gma.math.ufl.edu/exams/algebra-phd/}{Algebra PhD Exam}, September 2009\par}%
  \tagstructend%
  \vskip 0.35em%
  \tagmcbegin{artifact}%
    \hrule height 1pt%
  \tagmcend%
  \vskip 0.55em%
}
\makeatother
\title{Algebra PhD Exam, September 2009}
\author{}
\date{}
\begin{document}
\maketitle
\thispagestyle{empty}
\begin{examinstructions}
Do seven of the following eleven problems. Do not turn in more than seven problems. You must show your work. Answers with no work and/or no explanations will receive no credit. State clearly any theorem you use in your proofs.

\par
In the problems, \(\mathbb Z\), resp. \(\mathbb Q\), \(\mathbb C\), is the set of all integers, resp. of all rational numbers, of all complex numbers.
\end{examinstructions}
\begin{examproblems}{0}
\item
State and prove the Jordan–Hölder theorem for finite groups.
\item
This problem has two parts.
\begin{examsubparts}
\item[\textnormal{(a)}]
Show that the additive group of a field is never a free abelian group.
\item[\textnormal{(b)}]
Show that the multiplicative group of positive rational numbers is a free abelian group, and find a basis.
\end{examsubparts}
\item
A category~\(C\) is small if its class of objects is a set, and skeletal if for all~\(X\) and~\(Y\) in~\(C\), \(X\simeq Y\) implies \(X=Y\). Show that any small category is equivalent to a skeletal one (choose an object from each isomorphism class).
\item
Let~\(R\) be a commutative ring.
\begin{examsubparts}
\item[\textnormal{(a)}]
Define what it means for an~\(R\)-module~\(M\) to be flat.
\item[\textnormal{(b)}]
Show that the polynomial ring \(R[x]\), considered as an~\(R\)-module in the natural way, is flat.
\end{examsubparts}
\item
Show that if~\(R\) is a commutative ring with identity and every prime ideal of~\(R\) is finitely generated, then~\(R\) is Noetherian.
\item
Show that for any extension \(A\to B\) of commutative rings, any~\(A\)-module~\(M\) and any~\(B\)-module~\(N\), there is an isomorphism 
\[
\operatorname{Hom}_B(M\otimes_A B,N)\simeq \operatorname{Hom}_A(M,N_A)
\]
 functorial in~\(M\) and~\(N\) (here \(N_A\) denotes~\(N\) considered as an~\(A\)-module by means of the ring homomorphism \(A\to B\)).
\item
This problem has two parts.
\begin{examsubparts}
\item[\textnormal{(a)}]
Show that if \(L/F\), \(K/L\) are finite separable extensions, then so is \(K/F\).
\item[\textnormal{(b)}]
Do the same for purely inseparable extensions.
\end{examsubparts}
\item
This problem has three parts.
\begin{examsubparts}
\item[\textnormal{(a)}]
Show that if \(K/F\) is a normal extension and \(F\subseteq L\subseteq K\) is an intermediate field, then \(K/L\) is normal.
\item[\textnormal{(b)}]
Give an example where \(L/F\) is not normal.
\item[\textnormal{(c)}]
Give an example of normal extensions \(L/F\), \(K/L\) where \(K/F\) is not normal.
\end{examsubparts}
\item
Let~\(p\) be a prime and~\(K\) be the splitting field over \(\mathbb Q\) of \(X^p-1\).
\begin{examsubparts}
\item[\textnormal{(a)}]
Show that~\(K\) contains a unique quadratic extension of \(\mathbb Q\).
\item[\textnormal{(b)}]
When is this quadratic extension real? (You may assume as known the structure of the Galois group of \(K/\mathbb Q\). hint: if~\(\zeta\) is a primitive~\(p\)th root of unity, complex conjugation sends \(\zeta\mapsto\zeta^{-1}\)).
\end{examsubparts}
\item
Let~\(k\) be a field and~\(R\) a finite-dimensional~\(k\)-algebra. Show that if~\(R\) is semisimple as a ring, it is a direct sum of extension fields of~\(k\).
\item
Classify up to isomorphism all semisimple rings with \(720\) elements.
\end{examproblems}
\end{document}
